\begin{frame}{A universal quantum gate set}
  \small
  \begin{center}
    A \textbf{universal gate set} can approximate any desired quantum gate
    to any chosen accuracy.
  \end{center}

  \vspace{0.5em}
  \begin{columns}[T,onlytextwidth]
    \column{0.46\textwidth}
    \centering
    \textbf{Clifford gates}\par
    \vspace{0.35em}
    \begin{tikzpicture}[>=Stealth]
      \node[draw,rounded corners,fill=blue!10,minimum width=1.05cm,minimum height=0.7cm] (h) at (0,0) {$H$};
      \node[draw,rounded corners,fill=blue!10,minimum width=1.05cm,minimum height=0.7cm] (s) at (1.45,0) {$S$};
      \node[draw,rounded corners,fill=blue!10,minimum width=1.35cm,minimum height=0.7cm] (c) at (3.05,0) {CNOT};
    \end{tikzpicture}

    \vspace{0.4em}
    Hadamard makes superpositions; $S$ changes phase; CNOT links two qubits.
    Together, these are useful but \textbf{not universal}.

    \column{0.08\textwidth}
    \centering
    \vspace{1.3em}
    $+$

    \column{0.46\textwidth}
    \centering
    \textbf{Add a non-Clifford gate}\par
    \vspace{0.35em}
    \begin{tikzpicture}[baseline=(t.base)]
      \node[draw,rounded corners,fill=orange!18,minimum width=1.05cm,minimum height=0.7cm] (t) at (0,0) {$T$};
      \node[align=left] at (2.2,0) {$T=\begin{pmatrix}1&0\\0&e^{i\pi/4}\end{pmatrix}$};
    \end{tikzpicture}

    \vspace{0.4em}
    $T$ adds a $\pi/4$ phase to the $|1\rangle$ component. The combined set
    $\{H,S,T,\mathrm{CNOT}\}$ is called \textbf{Clifford+$T$}.
  \end{columns}

  \vspace{0.7em}
  \setbeamercolor{universaltakeaway}{fg=black,bg=blue!12}
  \noindent\makebox[\textwidth][c]{%
    \begin{beamercolorbox}[wd=0.96\textwidth,sep=0.65em,center,rounded=true]{universaltakeaway}
      \small\bfseries
      By combining these gates in sequences, we can approximate any quantum operation to arbitrary precision.
    \end{beamercolorbox}%
  }
\end{frame}
